\section{Introduction}
\label{sec:intro}

Suppose you need to test a C++ software stack that accepts inputs from multiple channels.
One component parses 3D binary assets, a second renders graphics and decodes audio, and a third accepts HTTP requests and file uploads over the network.
Developers commonly run coverage-guided fuzzers such as libFuzzer~\cite{serebryany2016libfuzzer} or AFL++~\cite{fioraldi2020aflplusplus} on individual functions.

These tools work well on standalone byte parsers.
They mutate raw input buffers and track branch coverage using compiler instrumentation.
However, these tools stop when software interacts with stateful runtime environments, external graphics displays, or asynchronous network event loops.
Traditional fuzzers either require writing bespoke mock harnesses for every component or suffer from low execution throughput when executing full integration builds.

We focus on the problem of building a unified adaptive fuzzing architecture that scales across diverse software layers.
Starting from a pure structured file format parser, our goal is to extend the same core mutator and scheduler to multimedia runtimes and asynchronous web frameworks without losing execution speed or bug-finding capability.

\begin{figure}[t]
\begin{lstlisting}[numbers=none,escapechar=@]
int HttpFileImpl::save(const std::string &p) const noexcept {
    if (!std::filesystem::exists(p)) // Throws on len > 255
        std::filesystem::create_directories(p);
} // Escaping exception in noexcept invokes std::terminate()
\end{lstlisting}
\caption{Motivating defect in \drogon (\textbf{D1}, simplified). Calling throwing \texttt{std::filesystem::exists} inside a \texttt{noexcept} method causes \texttt{std::terminate()} to abort the process when an upload filename exceeds 255 bytes.}
\label{fig:motivating}
\end{figure}

Consider the defect in Figure~\ref{fig:motivating}, discovered by our fuzzer in the \drogon C++ web framework.
The method \texttt{HttpFileImpl::save} was declared \texttt{noexcept}.
When an application saves an uploaded file using a client-controlled filename exceeding the filesystem limit (255 bytes on Linux), the call to \texttt{std::filesystem::exists} threw a \texttt{std::filesystem\_error}.
Because the enclosing method was marked \texttt{noexcept}, the C++ runtime intercepted the uncaught exception and immediately invoked \texttt{std::terminate()}.
A single malformed upload killed the entire multi-threaded web server process.
Across the three evaluated codebases, we uncovered 16 security and reliability defects spanning memory corruption, buffer overflows, and state desynchronization.

Prior work on structure-aware fuzzing handles complex file grammars through custom mutators~\cite{aschermann2019nautilus}.
However, existing grammar fuzzers do not provide headless shims for interactive game engines or asynchronous network protocols.
Conversely, network fuzzers such as AFLNet~\cite{pham2020aflnet} send packets over real operating system sockets, which restricts throughput to hundreds of executions per second.

We present \tool, an adaptive fuzzing architecture designed to test heterogeneous C and C++ software stacks.
\tool operates across four core components:
(i)~an adaptive mutation scheduler governed by a Discounted-UCB multi-armed bandit algorithm (Section~\ref{sec:bandit_scheduler});
(ii)~structured mutation extended from binary assets to multimedia rendering through headless shims (Section~\ref{sec:shims});
(iii)~in-memory TCP socket mocks for asynchronous network servers (Section~\ref{sec:shims}); and
(iv)~independent coverage accounting via GNU linker symbol wrapping (Section~\ref{sec:linker_wrapping}).

We evaluated \tool across three software projects: the \fastgltf 3D parser~\cite{fastgltf_repo}, the \raylib game engine~\cite{raylib_repo}, and the \drogon asynchronous web framework~\cite{drogon_repo}.
Across 59 target binaries, the tool executed 807{,}520{,}128 campaign iterations.
\tool discovered 16 security vulnerabilities and defects, resulting in 6 upstream-confirmed defects (two merged pull requests and four maintainer commits) alongside 10 unfiled medium- and high-severity findings.
Our results show that modular headless shims allow an in-memory fuzzer to reach up to 53{,}732 executions per second on network parsers while discovering real-world vulnerabilities.

\paragraph{Contributions.}
\begin{itemize}
  \item Cross-domain fuzzer adaptation across file formats, multimedia runtimes, and network services (Section~\ref{sec:problem}).
  \item The \tool architecture combining structure-aware mutation, Discounted-UCB scheduling, and linker wrapping (Section~\ref{sec:approach}).
  \item Domain-specific headless shims for software rasterizers, audio streams, and asynchronous event loops (Section~\ref{sec:approach}).
  \item Evaluation across three open-source systems over 807M executions documenting 16 defects (Section~\ref{sec:eval}), with open artifacts (Section~\ref{sec:data}).
\end{itemize}
